% ---
% title: A medical test
% tags: bayes, total-probability
% ---

\begin{problem}
	In a country 6\% of the people are infected. A new device diagnoses an
	infected person as sick with probability $0.96$, and it diagnoses a healthy
	person as sick with probability $0.05$.
	(a) Find the probability that the device diagnoses a randomly chosen person
	as sick.
	(b) The device diagnoses a person as sick. What is the probability that this
	person is really infected?
\end{problem}

\begin{solution}
	\idea Split the population into two groups: infected and healthy. Use the
	law of total probability for (a). Then use Bayes' formula for (b). We know
	the chance of a ``sick'' diagnosis in each group; Bayes' formula turns this
	around and gives the chance of infection for a person with a ``sick''
	diagnosis.

	Let $I$ be the event that the chosen person is infected, and $D$ the event
	that the device diagnoses the person as sick. The data are
	\[ P(I)=0.06,\qquad P(I^c)=1-0.06=0.94,\qquad P(D\mid I)=0.96,\qquad
		P(D\mid I^c)=0.05. \]
	\step{(a)} Every person is either infected or healthy, and not both. So $I$
	and $I^c$ split $\Om$ into two parts, and $D$ splits into the two mutually
	exclusive pieces $DI$ and $DI^c$. Add their probabilities, and write each one
	with the multiplication rule $P(DI)=P(D\mid I)P(I)$:
	\[ P(D)=P(D\mid I)P(I)+P(D\mid I^c)P(I^c)
		=0.96\cdot0.06+0.05\cdot0.94=0.0576+0.047=0.1046. \]
	So about 10.5\% of all people get a ``sick'' diagnosis.
	\step{(b)} We need $P(I\mid D)$. By the definition of conditional
	probability, $P(I\mid D)=P(ID)/P(D)$. The numerator is the first term of the
	sum in (a), $P(ID)=P(D\mid I)P(I)=0.0576$. The denominator is the whole
	sum from (a). This is Bayes' formula:
	\[ P(I\mid D)=\frac{P(D\mid I)P(I)}{P(D\mid I)P(I)+P(D\mid I^c)P(I^c)}
		=\frac{0.0576}{0.1046}\approx0.5507. \]
	\begin{answer}
		(a) $P(D)=0.1046$ \quad (b) $P(I\mid D)\approx0.5507$
	\end{answer}
	\begin{remark}
		\textbf{Common mistake:} to answer $0.96$ in (b). The number $0.96$ is
		$P(D\mid I)$, the chance of a ``sick'' diagnosis for an infected person.
		The question asks the other way around, $P(I\mid D)$. These are very
		different here: only about 55\% of the people with a ``sick'' diagnosis
		are really infected. The reason is that healthy people are the large
		majority (94\%). Even a small error rate of 5\% among them gives
		$0.047$, almost as much as the $0.0576$ that comes from the infected
		people.

		The structure is the same as in Exercise~\exref{early-warning}: a prior
		probability, and the chance of a signal in each of the two groups.
		Infection plays the role of the event of interest, and the diagnosis
		plays the role of the warning signal.
	\end{remark}
\end{solution}
